\renewcommand\S{\mathbb{S}}

\begin{problem}[1]
  Let $f : \S^n \to \S^n$ be a map with degree $\deg(f) \ne 1$. Show that there is a point $x \in \S^n$ with $f(x) = -x$.
\end{problem}

\begin{solution}
  Assume that $f(x) \ne -x$ for every $x \in \S^n$. Then, we can define a homotopy $H \times \S^n \times I \to \S^n$ by
  \[%
    H(x, t) = \frac{(1 - t)f(x) + tx}{\|(1 - t)f(x) + tx\|}
  .\]%
  Since the $f(x) \ne -x$, the denominator is never zero, so $H$ is well-defined. At $t = 0$, we get $H(x, 0) = f(x)$, and at $t = 1$, we get $H(x, 1) = x$. Therefore, $H$ is a homotopy between $f$ and $\Id_{\S^n}$, which is a contradiction, since maps that are homotopic have the same degree. Thus, there must be a point $x \in \S^n$ with $f(x) = -x$.
\end{solution}

\begin{problem}[2]
  Let $f : \S^n \to \S^n$ be a map with $n$ even. Show that there is a point $x \in \S^n$ with $f(f(x)) = x$.
\end{problem}

\begin{solution}
  Let $f : \S^{2k} \to \S^{2k}$ for some integer $k$. The only way for $f(x) \ne -x$, the $\deg f \ne -1$ (by the contrapositive of Corollary 9.3). Since $\deg f \circ f = (\deg f)^2 = 1$, then $\deg f = 1$. Therefore, $f \sim \Id_{\S^{2k}}$. Since $\Id_{\S^{2k}} \circ \Id_{\S^{2k}} = \Id_{\S^{2k}}$, then $f \circ f \sim \Id_{\S^{2k}}$. Thus, by Corollary 9.4, there exists a point $x \in \S^{2k}$ with $f(f(x)) = x$.
\end{solution}

\begin{problem}[3]
  Show that, if $n$ is odd, $-\Id_{\S^n} \sim \Id_{\S^n}$. (Hint: Try doing it when $n = 1$ first.)
\end{problem}

\begin{solution}
  Let $f(x) = \Id_{\S^n}$ and $g(x) = -\Id_{\S^n}$. We have $\deg f = 1$. Since $g(x)$ has no fixed points, then by Corollary 9.3, we have $\deg g = (-1)^{n + 1} = 1$ (since $n$ is odd). Thus, $f \sim g$, so $-\Id_{\S^n} \sim \Id_{\S^n}$.
\end{solution}

\begin{problem}[4]
  Let $K$ be a nonempty connected simplicial complex, and suppose that $f : |K| \to |K|$ is homotopic to a constant map. Show that there exists $x \in |K|$ with $f(x) = x$.
\end{problem}

\begin{solution}
  Since $f \sim g$, where $g(x) = x_0$, for every $x$, then $f_* = g_*$. On $H_0(K)$, we have $f_* = g_* = \Id_{H_0(K)}$, since $K$ is connected. Therefore, $\Lambda_f = 1$. By the Lefschetz fixed point theorem, there exists a point $x \in |K|$ with $f(x) = x$.
\end{solution}

\begin{problem}[5]
  Suppose that $f : |K| \to |K|$ is simplicial and that the set $S := \{x~|~f(x) = x\}$ is finite. You may assume that every element of $S$ is a vertex of $K^1$, which we stated but did not prove earlier in the course. More precisely, $S = \{\hat{A}~|~A \in K, f(A) = A\}$.
  \begin{enumerate}
    \item Show that, if $f(A) = A$, then there is no proper face $B \subsetneq A$ such that $f(B) = B$.
    \item Suppose that $f(A) = A$. Show that one can order the vertices of $A(v_0, \dots, v_i)$ such that $f(v_j) = v_{j+1}$ for all $0 \le j \le i - 1$ and $f(v_i) = v_0$.
    \item Let $\sigma = (v_0, \dots, v_i)$ be the orientation of $A$ that you constructed in part (ii). Show that $f_i(\sigma) = (-1)^i\sigma$.
    \item Show that $\Lambda_f = |S|$.
  \end{enumerate}
\end{problem}

\begin{solution}[(i)]
  Assume that there is a proper face $B \subsetneq A$ such that $f(B) = B$. Then, $\hat{B} \in S$, which contradicts the assumption that $S$ is finite. Therefore, there is no proper face $B \subsetneq A$ such that $f(B) = B$.
\end{solution}

\begin{solution}[(ii)]
  Let $A = (v_0, \dots, v_i)$ be a simplex such that $f(A) = A$. Since $f$ is simplicial, we have $f(v_j) \in A$ for all $0 \le j \le i$. Since there is no proper face $B \subsetneq A$ such that $f(B) = B$, then $f(v_j) \ne v_j$ for all $0 \le j \le i$. Therefore, we can order the vertices of $A$ such that $f(v_j) = v_{j+1}$ for all $0 \le j \le i - 1$ and $f(v_i) = v_0$.
\end{solution}

\begin{solution}[(iii)]
  Let $\sigma = (v_0, \dots, v_i)$ be the orientation of $A$ that we constructed in part (ii). Then, we have
  \[%
    f_i(\sigma) = (f(v_0), \dots, f(v_i)) = (v_1, \dots, v_i, v_0) = (-1)^i\sigma
  .\]%
  Therefore, $f_i(\sigma) = (-1)^i\sigma$.
\end{solution}

\begin{solution}[(iv)]
  Let $S = \{\hat{A}_1, \dots, \hat{A}_m\}$, where $A_j$ is a simplex such that $f(A_j) = A_j$ for all $1 \le j \le m$. Then, we have
  \[%
    \Lambda_f = \sum_{i = 0}^n (-1)^i\Tr(f_i) = \sum_{j = 1}^m (-1)^{\dim A_j}\Tr(f_{\dim A_j}) = \sum_{j = 1}^m (-1)^{\dim A_j}(-1)^{\dim A_j} = m
  .\]%
  Therefore, $\Lambda_f = |S|$.
\end{solution}
